\mychapter{Dataset}{Dataset}{}\label{chap:dataset}

\mysection{Background}{Background}\label{sec:background}

The dataset used contains three 72 hour recordings, recorded in False Bay, near Gordons Bay, between 22 April 2023 and 22 June 2023. 
The recordings were taken at a depth of 20 to 30 meters at the coordinates, S 34.16813deg, E18.758718deg. 

The device used for the recordings was a single omnidirectional hydrophone, designed and built by a Stellenbosch University research 
group from the Electrical and Electronic Engineering Department.

The recorded data has a sampling rate of 8kHz and a 16-bit resolution. The data was extracted as a 72 hour wav file from an SD card 
used by the recording device.\cite{data_set}

\mysection{Labeling}{Labeling}\label{sec:labeling}

Detection tables, containg start and end times of calls within the each file, were provided from a previous research project which used a 
band-limited energy detector to flag calls. This method unfortunately did not detect all the calls within the audio files and none of the 
calls were labelled as a specific type of call. 

Another detection table was procured. This table contained more detections for the first 72 hour file and contained call labels for all three files. 
This meant that there were sufficient labels for the first 72 hour recording of the dataset and there were partial labels for the second and third 72 
hour recordings within the dataset.\cite{data_set}

\mysection{Data Splitting}{Data Splitting}\label{sec:data-splitting}

For this project the data was plit in three sections. Each of the 72 hour audio files was chosen for a different section. Given that only 
the first of the three 72 hour audio files was well labeled, it was decided that this would be used for training and testing. The second 
audio file was used for hyperparameter tuning and the third was used to select templates.

\mysubsection{Template Selection}{Template Selection}\label{sec:temp_sel}

The audio file which was recorded starting on the 29th was used to select templates. These chosen templates were used for the comparison which 
determined whether an audio clip contained calls or not.

For the three different call type, single-tones, multi-tones and burst-tonals, the maximum number of templates per call type which could be 
extracted such that each call type had the same amount of templates was 10. The limiting call type was burst-tonals.

To ensure that the 10 templates chosen were representative of varying amounts of background noise templates were chosen by looking at spectrograms 
of detected calls throughout the audio file. The visual inspection was also necessary since some of the calls labeled as single-tones looked more like 
multi-tones and some of the multi-tones could actually be classified as burst-tonals.

\mysubsection{Hyperparameter Tuning Template Selection}{Hyperparameter Tuning Template Selection}\label{sec:tune_temp_sel}

The audio file which was recorded starting on the 26th was used to select calibration templates. These templates were used to determine below which threshold a 
segment of audio could be considered a call. For the calibration templates 10 templates were chosen for each call type. This was again limited by the burst-tonals 
maximum number of templates. In this file there were however only 9 labeled burst-tonals, but when selecting multi-tones it was found that a call originally labeled 
as a multi-tone could in fact be seen as a burst-tonal. 

Besides these 30 templates, 10 for each call type, 100 background templates were also selected. Within each audio file most of the time is made up of background 
noise and only some parts actually contain calls. This was simulated by selecting many more background templates that call-type templates. Since there is no table 
labeling background noise background segments of varying sizes were extracted from the file and 100 templates were manually selected. The reason for the manual process 
links with the unreliability of the labelled data. The audio file used was only labeled by the band-limited energy detector and thus did not label all the calls 
within the file. This means that some of the samples extracted as background audio actually contained calls that had never been labeled. Manually selecting each template 
ensured that all the templates chosen were in fact background noise only.

\mysubsection{Evaluation Data Selection}{Evaluation Data Selection}\label{sec:eval_dat_sel}

The entire first audio file which was recorded starting on the 22nd was used for the final evaluation of the whale call detection model. No templates were extracted, 
instead the entire file was sweeped to search for calls.